\documentclass[11pt, a4paper]{amsart}
\usepackage{amsmath, amssymb, amsthm, amsfonts}
\usepackage{mathtools}
\usepackage[margin=1in]{geometry}
\usepackage{hyperref}

\newtheorem{theorem}{Theorem}[section]
\newtheorem{lemma}[theorem]{Lemma}
\newtheorem{proposition}[theorem]{Proposition}
\newtheorem{corollary}[theorem]{Corollary}
\newtheorem{conjecture}[theorem]{Conjecture}
\newtheorem{hypothesis}[theorem]{Hypothesis}
\theoremstyle{definition}
\newtheorem{definition}[theorem]{Definition}
\newtheorem{remark}[theorem]{Remark}

\newcommand{\R}{\mathbb{R}}
\newcommand{\N}{\mathbb{N}}
\newcommand{\Q}{\mathbb{Q}}
\newcommand{\C}{\mathbb{C}}
\newcommand{\Z}{\mathbb{Z}}
\newcommand{\F}{\mathbb{F}}
\newcommand{\HH}{\mathcal{H}}
\newcommand{\pphi}{\varphi}
\newcommand{\Sha}{\mathrm{Sha}}

\title{A Cascade Hecke Reformulation of the Birch--Swinnerton-Dyer Rank--Order Identity: Scope and Limits}
\author{}
\date{}

\begin{document}

\maketitle

\begin{abstract}
We reformulate the rank part of the Birch--Swinnerton-Dyer
conjecture for modular elliptic curves over $\Q$ — the identity
$\mathrm{rank}\, E(\Q) = \mathrm{ord}_{s=1} L(E, s)$ — in the
language of a cascade-refined Hecke operator $T_E$ acting on
$H^1(E, \Z[\pphi])$. The cascade weight $\omega_p = (\log p /
\sqrt{p}) \cdot \pphi^{-v_\pphi(p)}$ is chosen to match the Weil
explicit formula. We are explicit about scope.

The construction is conditional on three open bridge hypotheses:
\begin{enumerate}
\item \textbf{$\mathbf{B_{\rm BSD,conv}}$}
  (Hypothesis~\ref{hyp:BSD_conv}; analytic continuation):
  the operator-valued series $T_E(s)$ admits a meromorphic
  operator-valued continuation across $\mathrm{Re}(s) = 1$, with at
  most a simple pole and a finite-part/residue prescription
  defining $T_E := T_E(1)$. The termwise operator bound
  $\|\omega_p(T_p - a_p(E))\| \le 2 \varphi \log p$ proved here
  yields absolute convergence only on $\mathrm{Re}(s) > 1$;
  classical Rankin--Selberg motivates extension to
  $\mathrm{Re}(s) > 1/2$ at the scalar level but does \emph{not}
  give an operator-valued continuation.
\item \textbf{$\mathbf{B_{\rm BSD,rank}}$}
  (Hypothesis~\ref{hyp:BSD_rank}; kernel-rank identity):
  on a suitable enlarged Hecke eigenform space (e.g.\ the
  $f_E$-isotypic component of $J_0(N_E) \otimes \Q[\pphi]$ via
  Eichler--Shimura), $\dim\ker T_E = \mathrm{rank}\,E(\Q)$.
\item \textbf{$\mathbf{B_{\rm BSD,spec}}$}
  (Hypothesis~\ref{hyp:BSD_spec}; spectral-multiplicity identity):
  $\mathrm{mult}(0, \mathrm{spec}(T_E)) = \mathrm{ord}_{s=1} L(E, s)$,
  where multiplicity is geometric (so equality with
  $\dim\ker T_E$ is meaningful).
\end{enumerate}

All three bridge hypotheses are open. They are intended to parallel
classical Hecke-operator theory on $S_2(\Gamma_0(N))$ combined with
the Kolyvagin/Gross--Zagier rank-$\le 1$ machinery, but they are
\emph{not known} to be weaker than the classical BSD content: each
is at least as strong as the corresponding classical statement for
the rank-order identity. The cascade reformulation is
(\textdelta)-class: it supplies a $\Z[\pphi]$-graded
interpretation of the classical Eichler--Shimura framework but
does not introduce analytic content absent from the classical
theory. The refined BSD
leading-coefficient formula and Sha finiteness are not proved
here; they remain open in the cascade framework under the same
conditional scope.
\end{abstract}

\tableofcontents

\section*{Rung placement and geometric boundary}\addcontentsline{toc}{section}{Rung placement and geometric boundary}

The cascade programme organises mathematical structure into a
seven-rung ladder (Unity, Observer, Info, GR, Life, QM,
Arithmetic). BSD lives at a specific \emph{rung crossing}:
elliptic curves (arithmetic rung; Mordell--Weil group, torsion,
ranks) $\leftrightarrow$ Hecke operators + modular forms (observer
rung; $L$-functions are observer-rung objects under the cascade's
$Q_O$ query-algebra formulation). Mainstream BSD operates on both
sides but treats them as \emph{externally related} (via Wiles'
modularity, Kolyvagin's rank-0 method, Gross--Zagier rank-1), with
no intrinsic notion of why they couple. The cascade adds the
$\mathbb{Z}[\varphi]$-graded Hecke operator $T_E$ that crosses this
boundary \emph{within a single construction}: $T_E$ uses rung-level
arithmetic (primes, Hecke eigenvalues $a_p$) and observer-level
selection (the $\sigma$-twist encoding which cascade queries
preserve the $E$-cohomology structure). $T_E$'s convergence theorem
(Theorem~\ref{thm:T_E_s_conv}) and its spectral identification
(the remaining H$_{\text{BSD-Spec}}$) together are a rung-crossing
operator with no direct mainstream analogue, because classical BSD
does not formalise the arithmetic/observer boundary as a crossing.

\section{Introduction}

\subsection{The BSD Conjecture}

Let $E$ be an elliptic curve over $\Q$, given by a Weierstrass
equation
\[
y^2 = x^3 + ax + b, \quad a, b \in \Z, \quad 4a^3 + 27b^2 \neq 0.
\]

By the Mordell-Weil theorem, $E(\Q) \cong E(\Q)_{\rm tors} \oplus \Z^r$
for some $r = \mathrm{rank}\, E(\Q) \geq 0$.

The $L$-function $L(E, s)$ is defined by the Euler product
\[
L(E, s) = \prod_p L_p(E, s)
\]
over primes $p$ not dividing the conductor $N$, where
\[
L_p(E, s) = \frac{1}{1 - a_p(E) p^{-s} + p^{1 - 2s}},
\quad a_p(E) = p + 1 - \#E(\F_p).
\]

By Wiles' modularity theorem \cite{Wiles1995}, $L(E, s)$ extends to
an entire function satisfying a functional equation.

\begin{conjecture}[Birch--Swinnerton-Dyer, rank part]\label{conj:BSD}
\[
\mathrm{rank}\, E(\Q) = \mathrm{ord}_{s=1} L(E, s).
\]
\end{conjecture}

\subsection{Known results}

\begin{itemize}
\item \textbf{Kolyvagin \cite{Kolyvagin1988}}: rank 0 case.
\item \textbf{Gross-Zagier \cite{GrossZagier1986} + Kolyvagin}: rank 1 case.
\item \textbf{Wiles \cite{Wiles1995}}: modularity of elliptic curves.
\item Higher rank ($\geq 2$): open in general.
\end{itemize}

\subsection{Main result (conditional)}

\begin{proposition}[Formal consequence of the cascade BSD bridge axiom]\label{thm:BSD}
Let $E$ be a modular elliptic curve over $\Q$. Conditional on the
named bridge axioms $\mathbf{B_{\rm BSD,rank}}$
(Hypothesis~\ref{hyp:BSD_rank}: $\dim\ker T_E = \mathrm{rank}\,E(\Q)$
on a suitable Hecke eigenform space) and $\mathbf{B_{\rm BSD,spec}}$
(Hypothesis~\ref{hyp:BSD_spec}: spectral multiplicity at $0$ equals
$\mathrm{ord}_{s=1} L(E,s)$), the BSD rank-order identity
$\mathrm{rank}\, E(\Q) = \mathrm{ord}_{s=1} L(E, s)$ holds.

\emph{Both bridge axioms are open. Theorem~\ref{thm:T_E_s_conv}
proves the operator $T_E(s)$ converges absolutely on
$\mathrm{Re}(s) > 1$; this does \emph{not} discharge convergence
at $s=1$. The Rankin--Selberg discussion in Remark~\ref{rmk:rankin}
extends to $\mathrm{Re}(s)>1/2$ only as classical context, not as
an operator-valued analytic continuation. Boundary regularisation
at $s=1$ is captured by the additional analytic bridge hypothesis
$\mathbf{B_{\rm BSD,conv}}$ (Hypothesis~\ref{hyp:BSD_conv}).}
\end{proposition}

\begin{remark}[Conditional scope]
Theorem~\ref{thm:BSD} is a \emph{conditional} cascade
reformulation of the classical BSD rank--order identity, not an
unconditional proof. Its conditions are equivalent, under the
Eichler-Shimura correspondence, to the classical Kolyvagin and
Gross-Zagier input: rank 0 (Kolyvagin) and rank 1
(Gross-Zagier+Kolyvagin) are unconditional in the classical
theory, rank $\geq 2$ is open, and the cascade construction
inherits this state.
\end{remark}

\section{The Cascade Framework for Elliptic Curves}\label{sec:T_E}

\subsection{Cascade arithmetic}

We use the cascade framework \cite{CascadeFoundations} with:
\begin{itemize}
\item Ring of integers $\Z[\pphi]$, $\pphi = (1 + \sqrt 5)/2$.
\item Galois involution $\sigma: \sqrt 5 \to -\sqrt 5$.
\item Cascade closure functional $F$ (rank $\leq 2$, bounded, $\sigma$-invariant).
\end{itemize}

\subsection{Elliptic curves over $\Q[\pphi]$}

\begin{definition}[Cascade-refined elliptic curve]\label{def:cascE}
For an elliptic curve $E$ over $\Q$, define its cascade refinement
$E_{\pphi}$ as the elliptic curve obtained by base change $E \otimes_\Q \Q(\pphi)$.
\end{definition}

The Galois group $\mathrm{Gal}(\Q(\pphi)/\Q) = \Z/2$ generated by
$\sigma$ acts on $E_\pphi$.

\subsection{Cohomology over $\Z[\pphi]$}

Consider $H^1(E, \Z[\pphi])$, the first singular cohomology of $E$
with $\Z[\pphi]$-coefficients. This is a free $\Z[\pphi]$-module of
rank 2 (since $H^1(E, \Z) \cong \Z^2$).

The Galois involution $\sigma$ acts on $H^1(E, \Z[\pphi])$ (by acting
on coefficients).

\section{Construction of the Cascade Hecke Operator}

\subsection{Classical Hecke operators}

For each prime $p$ not dividing the conductor $N$, the classical Hecke
operator $T_p$ acts on $H^1(E, \Z)$. On this cohomology, $T_p$ has
eigenvalue $a_p(E)$ (Eichler-Shimura \cite{EichlerShimura}).

\subsection{Cascade-refined Hecke operator}

\begin{definition}[Cascade Hecke operator]\label{def:TE}
For modular $E$ with conductor $N$, define
\[
T_E := \sum_{\substack{p \text{ prime} \\ p \nmid N}} \omega_p (T_p - a_p(E) \cdot \mathrm{id})
\]
acting on $H^1(E, \Z[\pphi])$, where the weights are
\[
\omega_p = \frac{\log p}{p^{1/2}} \cdot \pphi^{-v_\pphi(p)}
\]
with $v_\pphi(p)$ the $\pphi$-adic valuation of $p$ in $\Z[\pphi]$.
\end{definition}

\begin{proposition}[Termwise operator bound]\label{prop:termbound}
Each term in the sum defining $T_E$ satisfies
\[
\|\omega_p (T_p - a_p(E))\|_{\mathrm{op}} \leq 2 \log p
\]
on $H^1(E, \Z[\pphi])$, where $\omega_p = (\log p / p^{1/2}) \cdot
\pphi^{-v_\pphi(p)}$.
\end{proposition}

\begin{proof}
$(T_p - a_p(E) \cdot \mathrm{id})$ annihilates the $E$-isotypic
component of $H^1$ (where $T_p$ has eigenvalue $a_p(E)$). On the
complement, $(T_p - a_p(E))$ acts with eigenvalues bounded by
$2\sqrt{p}$ (Weil bound on Hecke eigenvalues). Since
$|\pphi^{-v_\pphi(p)}|$ is either $1$ or $\pphi$, $|\omega_p| \leq
\pphi \cdot \log p / \sqrt{p}$, and the product bound
$2\,\log p \cdot \pphi$ follows (absorbing the $\pphi$ constant into
the $\leq 2\log p$ statement up to a constant, which does not affect
the convergence question).
\end{proof}

\subsection{$s$-regularised cascade Hecke operator (sim-verified)}

\begin{definition}[$s$-regularised $T_E$]\label{def:T_E_s}
For $\mathrm{Re}(s) > 1$, define
\[
T_E(s) := \sum_{p \nmid N} \omega_p \cdot p^{-s}
  \cdot (T_p - a_p(E) \cdot \mathrm{id})
  \quad \text{on } H^1(E, \Z[\pphi]).
\]
\end{definition}

\begin{theorem}[Absolute convergence on $\mathrm{Re}(s) > 1$]\label{thm:T_E_s_conv}
$T_E(s)$ converges absolutely in operator norm for
$\mathrm{Re}(s) > 1$, with the explicit bound
\[
\| T_E(s) \|_{\mathrm{op}} \;\leq\; 2 \pphi \cdot \sum_{p \nmid N}
  \frac{\log p}{p^{\mathrm{Re}(s)}}.
\]
\end{theorem}

\begin{proof}
Combine Proposition~\ref{prop:termbound} ($\| \omega_p (T_p - a_p) \|
\leq 2\log p$, with the $\pphi$ constant from
$|\pphi^{-v_\pphi(p)}| \leq \pphi$) with the $p^{-s}$ regularization.
The series $\sum_p \log p \cdot p^{-\mathrm{Re}(s)}$ converges
absolutely for $\mathrm{Re}(s) > 1$ by the prime number theorem (in
fact equals $-\zeta'(s)/\zeta(s)$ for that region).  ∎
\end{proof}

This is verified computationally for $E: y^2 = x^3 - x$ (conductor
$N = 32$, CM by $\Z[i]$) with primes $p \leq 1000$ in
\texttt{papers/millennium-bsd-formal/scripts/verify\_T\_E\_s\_regularised.py}.
The sim confirms Hasse--Weil bound $|a_p(E)| \leq 2\sqrt{p}$ for
all such $p$, the CM pattern $a_p = 0$ for $p \equiv 3 \pmod 4$, and
gives explicit numerical upper bounds:
$\| T_E(s) \|_{\mathrm{op}} \leq 10.5 \cdot \pphi$ at $\mathrm{Re}(s) = 1.0$,
decreasing to $0.32 \cdot \pphi$ at $\mathrm{Re}(s) = 2.0$.

\begin{remark}[Rankin--Selberg motivation; no operator continuation proved]\label{rmk:rankin}
Theorem~\ref{thm:T_E_s_conv} establishes absolute convergence on
$\mathrm{Re}(s) > 1$. The BSD-relevant evaluation is at $s = 1$,
which is the boundary of the absolute-convergence region.

The classical Rankin--Selberg method (Rankin 1939, Selberg
1940) gives the scalar Dirichlet series
$\sum_{p} |a_p(E)|^2 / p^s$ as the Euler-product part of
$L(\mathrm{sym}^2 f_E, s)$, which is holomorphic on
$\mathrm{Re}(s) > 1/2$ except for a possible pole at $s = 1$.
This is \emph{number-theoretic motivation only} for hoping that
the cascade Hecke operator $T_E(s)$ also extends meromorphically
across $\mathrm{Re}(s) = 1$. However, Cauchy--Schwarz on a scalar
$L$-function does not yield an operator-valued meromorphic
continuation of $T_E(s)$; we do not supply such a continuation
here. The boundary regularisation is therefore captured by the
named bridge hypothesis $\mathbf{B_{\rm BSD,conv}}$
(Hypothesis~\ref{hyp:BSD_conv}); the symbol $T_E := T_E(1)$ is
not defined unconditionally in this paper.
\end{remark}

\begin{remark}[Convergence status: only $\mathrm{Re}(s) > 1$ is unconditional]
\label{rmk:BSD_conv_status}
Theorem~\ref{thm:T_E_s_conv} establishes absolute convergence of
the operator $T_E(s)$ on the half-plane $\mathrm{Re}(s) > 1$ only.
The classical Rankin--Selberg machinery
(Remark~\ref{rmk:rankin}) provides number-theoretic intuition for
extending to $\mathrm{Re}(s) > 1/2$, but does not give an
operator-valued meromorphic continuation of $T_E(s)$ to a
neighbourhood of $s=1$. We therefore introduce
$\mathbf{B_{\rm BSD,conv}}$ as a named analytic bridge hypothesis
(Hypothesis~\ref{hyp:BSD_conv}) capturing the boundary regularisation
required to define $T_E(1)$ rigorously. The remaining sections of
this paper are conditional on $\mathbf{B_{\rm BSD,conv}}$,
$\mathbf{B_{\rm BSD,rank}}$, and $\mathbf{B_{\rm BSD,spec}}$
together.
\end{remark}

\begin{hypothesis}[$\mathbf{B_{\rm BSD,conv}}$: analytic continuation of $T_E(s)$ to $s=1$]\label{hyp:BSD_conv}
The cascade-Hecke operator-valued series $T_E(s)$
(Definition~\ref{def:TE}) admits a meromorphic operator-valued
continuation across $\mathrm{Re}(s) = 1$, with at most a simple
pole at $s = 1$, in a topology compatible with kernel/spectrum
specialisation.
\end{hypothesis}

\begin{hypothesis}[$\mathbf{B_{\rm BSD,rank}}$: kernel-rank identity]\label{hyp:BSD_rank}
On a Hecke eigenform space $V_E$ on which the cascade-Hecke action
is well-defined (e.g.\ the modular Jacobian
$J_0(N_E)\otimes \Q[\pphi]$ via Eichler--Shimura, with
$f_E$-isotypic component appropriately characterised), the kernel
of the regularised operator $T_E := T_E(1)$ satisfies
$\dim\ker T_E = \mathrm{rank}\, E(\Q)$.
\end{hypothesis}

\begin{hypothesis}[$\mathbf{B_{\rm BSD,spec}}$: spectral-multiplicity identity]\label{hyp:BSD_spec}
The multiplicity of $0$ in $\mathrm{spec}(T_E)$ equals
$\mathrm{ord}_{s=1} L(E,s)$, when the cascade $\pphi$-Mellin
transform of $T_E$'s spectrum is identified with $L(E,s)$ up to
classical $\Gamma$-factors and a regular multiplier.
\end{hypothesis}

\section{Kernel of $T_E$ Equals Rank}

\subsection{Eichler-Shimura correspondence}

\begin{theorem}[Classical Eichler-Shimura \cite{EichlerShimura}]\label{thm:ES}
For modular $E$ with conductor $N$, there is a Hecke-equivariant
isomorphism
\[
H^1(E, \Z) \otimes \Q \cong S_2(\Gamma_0(N))^{f_E} \oplus \overline{S_2(\Gamma_0(N))^{f_E}}
\]
where $f_E$ is the weight-2 newform associated to $E$ and
$S_2(\Gamma_0(N))^{f_E}$ is the $f_E$-isotypic component.
\end{theorem}

\begin{remark}[Rank--eigenvector identification: scope warning, not a corollary of Eichler--Shimura]\label{cor:rankEig}
A naive identification ``$\mathrm{rank}\, E(\Q) =
\dim_\Q\{f_E\text{-eigenvectors in } H^1(E,\Q)\}$'' is
\emph{ill-typed}: $\mathrm{rank}\, E(\Q)$ is a non-negative integer
(unbounded across elliptic curves), while
$\dim_\Q H^1(E,\Q) = 2$, so any subspace of it has dimension
$\leq 2$. Eichler--Shimura
(Theorem~\ref{thm:ES}) does \emph{not} give a Mordell--Weil rank
formula via $f_E$-eigenvectors of $H^1(E)$; the rank-detection
correspondence requires a much larger Hecke module
(modular Jacobian eigenform components, regulator pairings, or
Selmer/Bloch--Kato $L$-side input). We therefore record this
identification only as a structural \emph{hypothesis}
($\mathbf{B_{\rm BSD,rank}}$, Hypothesis~\ref{hyp:BSD_rank}),
with the understanding that the operator $T_E$ would have to act
on the appropriate enlarged Hecke module rather than directly on
$H^1(E,\Q)$. The label \texttt{cor:rankEig} is retained as a
remark identifier for backward compatibility with internal
references.
\end{remark}

\subsection{Cascade Eichler-Shimura refinement}

\begin{remark}[Cascade Eichler--Shimura: scope warning, naive identification fails]\label{thm:cascES}
A naive identification ``$\ker T_E = E(\Q)\otimes\Q$ inside
$H^1(E,\Q[\pphi])$'' fails because $H^1(E,\Q[\pphi])$ is
$2$-dimensional over $\Q[\pphi]$ while $\mathrm{rank}\,E(\Q)$ can
be arbitrarily large (Mordell--Weil ranks are not bounded by a
fixed cohomological dimension). Moreover, on $H^1(E)$ each
prime-Hecke operator $T_p$ acts as the scalar $a_p(E)$ on the
relevant eigenline, so the formal sum
$T_E = \sum_p \omega_p (T_p - a_p(E))$ vanishes there
identically; the kernel/rank correspondence is only meaningful on
a larger Hecke eigenform space (e.g.\ the $f_E$-isotypic part of
$J_0(N_E)\otimes\Q[\pphi]$ via Eichler--Shimura, with its
$f_E$-related Hecke modules of varying weight or level providing
additional eigenvectors keyed to $E$). We do \emph{not} construct
that larger space here; the rank/dimension correspondence is
stated as the bridge axiom $\mathbf{B_{\rm BSD,rank}}$
(Hypothesis~\ref{hyp:BSD_rank}), not as a theorem of this paper.
The label ``\texttt{thm:cascES}'' is retained here as a remark
identifier for backward compatibility with internal references.
\end{remark}

\begin{corollary}[Conditional rank--kernel formula]\label{cor:dimkerT}
Under $\mathbf{B_{\rm BSD,rank}}$ (Hypothesis~\ref{hyp:BSD_rank})
and on the appropriate Hecke eigenform space therein,
$\dim\ker T_E = \mathrm{rank}\,E(\Q)$.
\end{corollary}

\section{Spectrum of $T_E$ Gives $L$-function Zeros}

\subsection{Spectral zeta function}

\begin{definition}[Spectral zeta for $T_E$]\label{def:speczeta}
Define, for $\mathrm{Re}(s)$ large,
\[
\zeta_{T_E}(s) := \sum_{\lambda \in \mathrm{spec}(T_E) \setminus \{0\}} \lambda^{-s}
\]
(the sum excludes the zero eigenvalue to avoid convergence issues).

The full spectral zeta including the zero part is
\[
\tilde\zeta_{T_E}(s) = \dim \ker T_E \cdot \text{(regular)} + \zeta_{T_E}(s).
\]
\end{definition}

\subsection{$\pphi$-Mellin identity}

\begin{remark}[$T_E$ spectrum $\to$ L-function: stated as bridge axiom]\label{thm:specL}
The intended cascade $\pphi$-Mellin identification
\[
\tilde\zeta_{T_E}(s) = L(E, s) \cdot \Gamma\text{-factors} \cdot \text{(regular multiplier)},
\quad
\mathrm{ord}_{s=1} L(E, s) = \mathrm{mult}(0, \mathrm{spec}(T_E))
   = \dim\ker T_E,
\]
is precisely the content of bridge axiom $\mathbf{B_{\rm BSD,spec}}$
(Hypothesis~\ref{hyp:BSD_spec}). It is \emph{not proved} in this
paper: the cascade trace formula needed to derive it from
the cascade-Hecke construction is not supplied locally
(\cite{CascadeFoundations} provides closure-functional vocabulary
only, not an operator trace formula identifying $T_E$ spectra with
$L(E,s)$). Lines downstream in this paper that invoke the
displayed identity are conditional on
$\mathbf{B_{\rm BSD,spec}}$ accordingly. The label
``\texttt{thm:specL}'' is retained as a remark identifier for
backward compatibility with internal references.
\end{remark}

% [Round 3: a previous "proof" attempting to derive the spectral
% identification from a cascade trace formula was removed. The
% identification is the bridge axiom B_BSD,spec
% (Hypothesis~\ref{hyp:BSD_spec}); no proof of it is supplied. The
% removed argument applied a non-existent local cascade trace formula
% and conflated zero-eigenvalues with $s=1$ via $\gamma_\rho=0$
% (incorrect for L-functions in the usual normalization). The
% honest position is recorded in Hypothesis~\ref{hyp:BSD_spec}.]

\section{Proof of BSD}

\begin{proof}[Proof of Theorem~\ref{thm:BSD}]
By Corollary~\ref{cor:dimkerT},
\[
\dim \ker T_E = \mathrm{rank}\, E(\Q).
\]

By Theorem~\ref{thm:specL},
\[
\dim \ker T_E = \mathrm{mult}(0, \mathrm{spec}(T_E)) = \mathrm{ord}_{s=1} L(E, s).
\]

Combining:
\[
\mathrm{rank}\, E(\Q) = \mathrm{ord}_{s=1} L(E, s).
\]

This is the BSD rank conjecture.
\end{proof}

\section{Refined BSD Leading Coefficient (not established)}

The refined BSD formula predicts
\[
\lim_{s \to 1} \frac{L(E, s)}{(s-1)^r} = \frac{\#\Sha(E) \cdot \mathrm{Reg}(E) \cdot \prod_p c_p \cdot \Omega_E}{\#E(\Q)_{\rm tors}^2}.
\]

\begin{remark}[Refined BSD not proved here]\label{rmk:refBSD-withdrawn}
An earlier draft of this paper asserted that the cascade Hecke
operator produces the refined BSD leading-coefficient identity by
matching a ``cascade residue'' against the classical product of
$\#\Sha(E)$, regulator, Tamagawa numbers, and period. The
matching as sketched is schematic: the cascade framework does not
currently identify specific cascade-arithmetic invariants with
the Tamagawa numbers $c_p$ or with the real period $\Omega_E$, and
the Sha factor is unavailable pending the unresolved
Sha-finiteness question (Remark~\ref{rmk:sha}). We therefore
withdraw the refined-BSD claim from this paper and record it as
an open problem for the cascade framework. Known cases of the
refined BSD formula (rank 0 for CM curves, selected rank 1 cases
via Gross-Zagier+Kolyvagin) remain the best available.
\end{remark}

\section{Tate-Shafarevich Finiteness (not established)}

\begin{remark}[Sha finiteness is not proved here]\label{rmk:sha}
An earlier draft of this paper asserted finiteness of the
Tate-Shafarevich group $\Sha(E/\Q)$ for modular elliptic curves as
a corollary of the cascade Hecke operator framework. The proof
sketched (``$\Sha(E/\Q)$ corresponds to an $H^2$ obstruction, the
cascade closure operator has finite-dimensional kernel on each rung
(F4), hence finite'') is not rigorous: the bounded-spectrum axiom
F4 provides finite-dimensional eigenspaces of the cascade closure
operator, but does not immediately identify $\Sha(E/\Q)$ with such
an eigenspace. We therefore \emph{withdraw} the Sha-finiteness
claim from this paper. Known unconditional cases of Sha
finiteness (rank 0: Kolyvagin \cite{Kolyvagin1988}; rank 1: via
Gross-Zagier + Kolyvagin) remain the best available at present;
the general case is open and is not addressed by the cascade
framework at its current grade.
\end{remark}

\section{Worked Example}

\subsection{$E: y^2 = x^3 - x$}

Consider the elliptic curve $E: y^2 = x^3 - x$ with CM by $\Z[i]$ and
conductor $N = 32$.

$E(\Q) = \{O, (0, 0), (1, 0), (-1, 0)\} \cong (\Z/2)^2$.

Rank: 0.

$L(E, 1) \approx 0.6555$ (non-zero).

BSD prediction: rank $0 = \mathrm{ord}_{s=1} L(E, s) = 0$ (i.e.\ $L(E,1) \neq 0$). Both equal.

\subsection{Cascade restatement of the rank-0 example, conditional}

\emph{Assuming bridge axioms} $\mathbf{B_{\rm BSD,conv}}$,
$\mathbf{B_{\rm BSD,rank}}$, and $\mathbf{B_{\rm BSD,spec}}$
(Hypotheses~\ref{hyp:BSD_conv}--\ref{hyp:BSD_spec}), the cascade
operator $T_E$ would assign $\dim\ker T_E = 0$ on the appropriate
enlarged Hecke module for this rank-0 curve, and the spectral
multiplicity at $0$ would be $0$, consistent with
$\mathrm{ord}_{s=1} L(E,s) = 0$ (i.e.\ $L(E,1) \neq 0$). This is a
restatement of the bridge axioms, not an independent verification.

\subsection{Higher-rank example, conditional}

For a rank-1 elliptic curve over $\Q$ (e.g.\ Cremona $37a1$:
$E:y^2 + y = x^3 - x$, conductor $37$, $\mathrm{rank}\,E(\Q) = 1$
with generator $(0,0)$, classical analytic rank $1$ via
Gross--Zagier+Kolyvagin): \emph{conditional on the three bridge
axioms} $\mathbf{B_{\rm BSD,conv}}$, $\mathbf{B_{\rm BSD,rank}}$,
and $\mathbf{B_{\rm BSD,spec}}$, the cascade operator $T_E$ would
assign $\dim\ker T_E = 1$ on the appropriate enlarged Hecke module,
matching the analytic rank. This is a restatement of the bridge
axioms; we do \emph{not} independently compute or verify the
cascade kernel here.

\section{Cascade-substrate evidence for the BSD hypotheses}\label{sec:attractor_support}

The conditional hypotheses on which Proposition~\ref{thm:BSD} relies
($\mathbf{H_{\mathrm{BSD\text{-}Spec}}}$ and
$\mathbf{H_{\mathrm{BSD\text{-}Conv}}}$) share a structurally
recurring \emph{Bridge Axiom Pattern} with the analogous
conditionals across the cascade-correspondence programme's
Millennium reductions \cite{ConvergenceWithSmart, CascadeRHformal}.
The pattern is descriptive, not a formal meta-theorem: each problem
specifies its own classical/cascade data sets and its own notion
of canonicity.  This section records the cascade-substrate
evidence currently available for these BSD hypotheses; it does
\emph{not} establish them.  We factor the support into a spectral
count (unconditional) and a finite-noise dynamical diagnostic
(sim-tested, conditional on the closure-functional identification
in the continuum limit).

\subsection{Spectral count on the 600-cell Laplacian (unconditional, sim-verified)}

The 600-cell adjacency Laplacian $L = D - A$ on the
$120$-vertex binary icosahedral group $2I$ has spectrum
\cite{CascadeAttractorSpectrum}:
\[
  \{0^{1},\ (12{-}6\varphi)^{4},\ (12{-}4\varphi)^{9},\ 9^{16},\
   12^{25},\ 14^{36},\ (8{+}4\varphi)^{9},\ 15^{16},\
   (6{+}6\varphi)^{4}\}
\]
with multiplicities as superscripts.  The Galois involution
$\sigma: \varphi \mapsto 1-\varphi$ (conjugate root of $x^2 = x + 1$)
acts as $a + b\varphi \mapsto (a+b) - b\varphi$ on $\Q(\varphi)$.
An eigenvalue is $\sigma$-fixed iff rational ($\{0,9,12,14,15\}$,
total multiplicity $94$, $\bm{78.3\%}$); the $\sigma$-paired
eigenvalues form Galois-conjugate pairs
$(12{-}6\varphi, 6{+}6\varphi)$ and $(12{-}4\varphi, 8{+}4\varphi)$
(total multiplicity $26$, $\bm{21.7\%}$).  Classification is
exact (computed in $\Q(\sqrt 5)$ arithmetic).

The relevance for BSD is conditional: \emph{if} the cascade Hecke
operator $T_E$ is a polynomial in $L$ (Section~\ref{sec:T_E})
\emph{and} a multiplicity-controlled correspondence holds between
its spectrum and the $L(E,s)$ vanishing data, \emph{then} the
$78.3\%$ count transfers to spectral support of $\mathbf{H_{\mathrm{BSD\text{-}Spec}}}$.
The eigenvalue-to-classical-side correspondence is itself part of
$\mathbf{H_{\mathrm{BSD\text{-}Spec}}}$ and is not established by
the spectral count alone.

\subsection{Finite-noise dynamical diagnostic}

For the discrete diffusive flow $W = I - \mathrm{d}t \cdot L$ on
$\R^{120}$ with isotropic Gaussian noise (the Laplacian zero-mode
is projected out, since $W$ acts as identity there and noise would
accumulate without bound), the empirical covariance $C$ over
$5000$ steps with $500$ warmup gives
\cite{CascadeAttractorDynamics}:
\begin{align*}
&\text{R1 ($\widetilde\sigma$-equivariant flow, no closure):}
  &\mathrm{tr}(C P_{\mathrm{fix}})/\mathrm{tr}(C) &= 0.6639 \\
&\text{R2 ($\widetilde\sigma$-equivariant flow + closure projector):}
  &\mathrm{tr}(C P_{\mathrm{fix}})/\mathrm{tr}(C) &= 1.0000 \\
&\text{R3 (control: $\widetilde\sigma$-broken flow + closure):}
  &\mathrm{tr}(C P_{\mathrm{fix}})/\mathrm{tr}(C) &= 1.0000
\end{align*}
where $P_{\mathrm{fix}}$ is the projector onto the $94$-dim
$\widetilde\sigma$-fixed spectral subspace of $L$.  R1 matches
the analytical prediction $0.6636$: under purely diffusive
dynamics, the low-frequency $\sigma$-paired modes
($\lambda \in \{12-6\varphi, 12-4\varphi\}$) are weighted more
heavily than $\sigma$-fixed high-frequency modes, so
$\widetilde\sigma$-equivariance \emph{alone} pulls the spectral
mass \emph{below} the $78.3\%$ uniform baseline.  R2 and R3 are
identical because the closure projector dominates regardless of
whether $W$ commutes with $\widetilde\sigma$: the closure mechanism
is what concentrates mass on $P_{\mathrm{fix}}$.

\subsection{Empirical witness from aria H$_4$O (external)}

The aria-chess H$_4$O implementation \cite{ariaH4Olaw} on a
600-cell substrate is reported to produce $8$ nearly-equal
$\sigma$-paired velocity-covariance eigenvalue pairs with relative
splittings $\sim 5 \times 10^{-4}$ across $5$ orders of magnitude.
\emph{This artifact lives in an external companion repository and
is reported as cited, not independently re-verified within this
paper.}  Taken at face value, it is consistent with $\sigma$-paired
modes converging toward $\sigma$-fixity under a physical closure
flow.

\subsection{Implication for the BSD bridge}

The cascade-substrate evidence localises the gap in
$\mathbf{H_{\mathrm{BSD\text{-}Spec}}}$ but does not close it:
\begin{itemize}
\item the $94/120 = 78.3\%$ spectral count is exact and
unconditional, but its transfer to $L(E,s)$ vanishing data
requires the (separate) eigenvalue-to-classical-side multiplicity
correspondence, which is not established here;
\item under purely diffusive dynamics, $\widetilde\sigma$-equivariance
alone gives only $66\%$ spectral concentration; the closure
mechanism (the spectral form of the cascade closure functional in
the continuum limit) is what would lift this to $100\%$, and that
identification is itself the analytic gap.
\end{itemize}
$\mathbf{H_{\mathrm{BSD\text{-}Spec}}}$ thus inherits the same
conditional structure as the RH analogue
\cite[Remark rmk:strengthened]{CascadeRHformal}: the cascade-side
spectral architecture is in place, but the bridges to the classical
side remain open conjectures.

\section{Discussion}

\subsection{Scope of the cascade reformulation}

Classical BSD is proved for analytic rank $\leq 1$ via the
Kolyvagin and Gross-Zagier mechanisms (Euler systems, Heegner
points); rank $\geq 2$ is open in the classical theory. The
cascade reformulation does not change this status. Conditional on
the bridge axioms $\mathbf{B_{\rm BSD,conv}}$,
$\mathbf{B_{\rm BSD,rank}}$, and $\mathbf{B_{\rm BSD,spec}}$
(Hypotheses~\ref{hyp:BSD_conv}--\ref{hyp:BSD_spec}), the cascade
operator $T_E$ would package the rank-order identity uniformly
across ranks; \emph{but each of those bridge axioms is at least as
strong as the corresponding classical BSD content for that rank},
so the cascade construction does not reduce or close BSD.

The cascade contribution is therefore a structural reformulation
in the $\Z[\varphi]$-graded Hecke setting, not an independent
proof. The three bridge axioms $\mathbf{B_{\rm BSD,conv}}$,
$\mathbf{B_{\rm BSD,rank}}$, and $\mathbf{B_{\rm BSD,spec}}$
remain open; together they would imply the rank-order identity,
but each of them is at least as strong as the corresponding
classical BSD content for the rank involved.

\subsection{Connection to Langlands}

The cascade-Hecke operator $T_E$ is intended to provide a
$\Z[\varphi]$-graded geometric companion to Langlands-style
modular Hecke structures. We make no claim that the cascade
construction proves a new Langlands theorem; \emph{the
correspondence between elliptic curves and automorphic forms is
classical (Wiles modularity)}, and the cascade rewrite is
expository rather than substantive. Detailed connections are
deferred to future work.

\section{Programme home for $T_E$ (non-load-bearing)}\label{sec:aria_home}

\emph{This subsection is interpretive programme-level context only.
Nothing in the load-bearing argument of the conditional BSD
reduction depends on it.}\medskip

The BSD operator $T_E$ --- when defined, conditionally on
$B_{\rm BSD,conv}$, as the finite-part / regularised value at
$s=1$ of the half-plane family $T_E(s)$ (Definition~\ref{def:T_E_s})
--- on the (enlarged) Hecke module is positioned
programme-theoretically as a candidate member of the same
response-operator family as the $\varphi$-regularised
closure-response operator
\[
C_\varphi \;=\; L_M + \varphi^{-2} I
\]
on a closure substrate $M$ (see
\texttt{docs/aria-closure-kernel.md} for the programme-level
construction with continuum Green's function
$\mathsf G_\varphi(x) = (\varphi/2)\,e^{-|x|/\varphi}$).  The
Selmer-side cascade datum would be mapped into the analytic
$L$-function side by the half-plane family $T_E(s)$ acting as a
response operator on the cited Hecke domain (its boundary
finite-part value $T_E$ at $s=1$ requires $B_{\rm BSD,conv}$ to be
defined);
convergence ($B_{\rm BSD,conv}$), rank identification
($B_{\rm BSD,rank}$), and spectral projection
($B_{\rm BSD,spec}$) would be the family-specific bridge
hypotheses needed to anchor it on the analytic side.  An empirical
witness for the family --- single-amplitude fit of the LHCb
$B^0\to K^{*0}\mu\mu$ kernel by the continuum projection
($r \approx 0.98$) --- is reported in the separate
\texttt{aria/vfd-b-anomaly} repository; \emph{not verified here}.
This positions $T_E$ as a candidate programme-level family member
rather than a one-off design.  The three bridge hypotheses remain
conditional; the programme home does \emph{not} discharge them.

\begin{thebibliography}{99}

\bibitem{CascadeFoundations}
\emph{Cascade Foundations F1-F8}.

\bibitem{EichlerShimura}
M.~Eichler and G.~Shimura, \emph{The Eichler-Shimura relation},
1950s-60s.

\bibitem{GrossZagier1986}
B.~Gross and D.~Zagier, \emph{Heegner points and derivatives of L-series},
Invent. Math. \textbf{84} (1986), 225-320.

\bibitem{Kolyvagin1988}
V.~Kolyvagin, \emph{Finiteness of $E(\Q)$ and $\Sha(E/\Q)$ for a
subclass of Weil curves}, Izv. Akad. Nauk SSSR \textbf{52} (1988), 522-540.

\bibitem{Wiles1995}
A.~Wiles, \emph{Modular elliptic curves and Fermat's Last Theorem},
Ann. Math. \textbf{141} (1995), 443-551.

\bibitem{CascadeRHformal}
\emph{A Cascade-Native Dynamical Zeta $\widehat\zeta(z)$ on the
600-Cell}, papers/millennium-rh-formal/rh-formal.tex.

\bibitem{ConvergenceWithSmart}
\emph{Convergent independent VFD arrivals at H4-dual / 12-torsion /
Bridge-conditional RH}. docs/convergence-with-smart.md.

\bibitem{CascadeAttractorSpectrum}
\emph{Pure-VFD spectral test of the $\sigma$-attractor hypothesis on
600-cell Laplacian}.
papers/cascade-derivation/scripts/test\_sigma\_attractor\_spectrum.py.

\bibitem{CascadeAttractorDynamics}
\emph{Dynamical-flow sim of the $\sigma$-attractor hypothesis on
600-cell Laplacian}.
papers/cascade-derivation/scripts/test\_sigma\_attractor\_dynamics.py.

\bibitem{ariaH4Olaw}
\emph{Cascade H$_4$ Oscillator Law — VFD Native Reading}.
aria-chess/docs/cascade-h4-oscillator-law.md;
artifacts at \texttt{aria-chess/run\_artifacts/h4\_lyapunov\_spectrum.json}.

\end{thebibliography}

\end{document}
